\section{The Fourth Week and the Seven Great Antiphons}\label{sec:week-4}

The Fourth Sunday of Advent always falls between 18 and 24 December, so its
weekdays belong entirely to the final stretch before Christmas. In the shortest
Advents---when the First Sunday falls on 3 December, as in 2028 and 2034---the
Fourth Sunday falls on 24 December itself, and under the 1960 rubrics (RG
n.~30~a) the Vigil of the Nativity takes its place outright with no
commemoration of the Sunday. In every other year, the Roman church assembles on
this Sunday at the Basilica of the Twelve Holy Apostles, singing once more
\textit{Rorate caeli}, ``Drop down dew, ye heavens, from above, and let the
clouds rain the just'' (Is 45:8).

That ancient Advent cry holds a quiet image that the Fathers loved to pair with
Psalm 71~[72]:6, ``He shall come down like rain upon the fleece; and as showers
falling gently upon the earth.'' St.~Cyril of Jerusalem contrasts Christ's
first coming---``the unobserved, like rain on a fleece''---with his open return
in glory (Catechesis XV.1). St.~Bernard of Clairvaux dwells on Gideon's fleece
in his Second Homily on the \textit{Missus est}:

\begin{quote}
Gently at first, and without noise of human operation, He fell softly into the
Virgin's womb; afterwards, when the Apostles announced Him, it was with the
noise of words and the display of miracles.
\end{quote}

On the Octave Day of Christmas, the Church sings that exact verse at Lauds:
\textit{Sicut pluvia in vellus descendisti}, ``like rain upon the fleece.''

\subsection{The Fourth Sunday: Salvation Named Inside History}

The Collect of the Fourth Sunday repeats the opening cry of the First Sunday
with a sharpened urgency, admitting frankly that our own sins hold back what
God is ready to give:

\begin{massbox}{Collect of the Fourth Sunday of Advent}
\textbf{Latin.} Excita, quaesumus, Domine, potentiam tuam, et veni: et magna
nobis virtute succurre; ut per auxilium gratiae tuae, quod nostra peccata
praepediunt, indulgentia tuae propitiationis acceleret: Qui vivis et regnas cum
Deo Patre in unitate.

\textbf{English.} Bestir, O Lord, Thy might, we beseech thee, and come; and
with great power come to our aid, that, by the help of Thy grace, that which is
hindered by our sins may be hastened by Thy merciful forgiveness. Who livest.
\end{massbox}

In the Epistle (1 Corinthians 4:1--5), St.~Paul refuses to rush the verdict on
himself or anyone else: ``Therefore, judge not before the time: until the Lord
come, who both will bring to light the hidden things of darkness and will make
manifest the counsels of the hearts.'' Then the Gospel (Luke 3:1--6) repeats
the passage first heard on Ember Saturday, and its opening lines (Luke 3:1--2)
stop an attentive listener cold:

\begin{quote}
Now in the fifteenth year of the reign of Tiberius Caesar, Pontius Pilate being
governor of Judea, and Herod being tetrarch of Galilee, and Philip his brother
tetrarch of Iturea and the country of Trachonitis, and Lysanias tetrarch of
Abilina: Under the high priests Annas and Caiphas: the word of the Lord was
made unto John, the son of Zachary, in the desert.
\end{quote}

The Gospel of the Sunday right before Bethlehem refuses to drift into timeless
legend. It pins the word of God to a numbered year of a Roman emperor and names
the men who held civil and religious office in Palestine---above all Pontius
Pilate, Herod, Annas, and Caiphas. The very magistrate who will wash his hands
and hand Jesus over to the cross already stands in the doorway of the Fourth
Sunday of Advent. And into that crowded political map steps one unarmed man in
the desert crying from Isaiah 40:3--5: ``Prepare ye the way of the Lord, make
straight his paths. Every valley shall be filled and every mountain and hill
shall be brought low: and the crooked shall be made straight, and the rough
ways plain. And all flesh shall see the salvation of God.''

St.~Gregory the Great, whose Twentieth Homily is read at Matins on this Sunday,
refuses to let those images stay geographical: ``But the crooked places are
become straight, when the hearts of the wicked, perverted by a course of
injustice, are directed to the rule of justice. But the rough ways are changed
to smooth, when fierce and savage dispositions by the influence of Divine grace
return to gentleness and meekness.'' The Offertory turns quietly to Gabriel's
greeting (\textit{Ave, Maria}, Luke 1:28), and the Communion echoes Isaiah
7:14, \textit{Ecce, Virgo concipiet}, ``Behold a virgin shall conceive, and
bear a son.''

\subsection{The Seven Great O Antiphons (17--23 December)}

From 17 December through 23 December inclusive, the weekdays of Advent rise to
second-class rank (RG n.~24~a), and every evening at Vespers the Church frames
Mary's \textit{Magnificat} with one of the seven Great Antiphons, sung whole
both before and after the canticle. Because every one opens with the vocative
``O,'' generations of Christians have called them the Great O's.

They are among the oldest surviving treasures of Western evening prayer. The
late tenth-century Hartker Antiphonary of St.~Gall already copies these exact
seven antiphons in their exact Roman order---opening a series of twelve undated
Vespers antiphons for the week before Christmas---with \textit{O Sapientia}
first and \textit{O Emmanuel} seventh.

One calendar fact is easily misstated and worth getting right: the seven Great
Antiphons do \textit{not} all fall on weekdays. Because the Fourth Sunday of
Advent always falls between 18 and 24 December, at least one Great Antiphon
falls on Sunday in six out of seven calendar shapes; and in the shortest
Advents, when the First Sunday falls on 3 December (as in 2028 and 2034), the
Third Sunday itself falls on 17 December and takes \textit{O Sapientia} at
Vespers while 24 December becomes the Vigil. On any Sunday from 17 to 23
December, the Great Antiphon simply displaces the ordinary Sunday
\textit{Magnificat} antiphon; on a feast such as St.~Thomas when 21 December
falls on a weekday, it is kept as the commemoration of the Advent feria (when
21 December is the Fourth Sunday, the Sunday takes precedence and sings
\textit{O Oriens} directly).

Each antiphon addresses Christ by a messianic title drawn from the Old
Testament, praises what he has done or promised, and ends with a single sharp
cry---``Come!''---matched to that title. Together they read like a seven-day
litany gathering the whole longing of Israel before the door of Bethlehem:

{\footnotesize
\begin{longtable}{>{\raggedright\arraybackslash}p{0.10\linewidth}>{\raggedright\arraybackslash}p{0.20\linewidth}>{\raggedright\arraybackslash}p{0.22\linewidth}>{\raggedright\arraybackslash}p{0.40\linewidth}}
\toprule
\textbf{Date} & \textbf{Antiphon} & \textbf{Scriptural Root} & \textbf{Cry to the Coming Lord}\\
\midrule
\endhead
17 Dec & \textit{O Sapientia} (O Wisdom) & Ecclus 24:5; Wis 8:1 & Reaches from end to end mightily and sweetly: ``come, to teach us the way of prudence!''\\
18 Dec & \textit{O Adonai} (O Lord) & Ex 3:2, 6:6, 19--20 & Appeared to Moses in the burning bush and gave the law on Sinai: ``come, to redeem us with an outstretched arm!''\\
19 Dec & \textit{O Radix Jesse} (O Root of Jesse) & Is 11:10, 52:15 & Stands as an ensign of the peoples, sought by the Gentiles: ``come, to deliver us, make no tarrying!''\\
20 Dec & \textit{O Clavis David} (O Key of David) & Is 22:22, 42:7; Apoc 3:7 & Opens and none shuts, shuts and none opens: ``come, to bring out the prisoners from the prison, and them that sit in darkness, and in the shadow of death!''\\
21 Dec & \textit{O Oriens} (O Day-Spring) & Luke 1:78--79; Wis 7:26; Mal 4:2 & Brightness of everlasting light and Sun of justice: ``come, to give light to them that sit in darkness, and in the shadow of death!''\\
22 Dec & \textit{O Rex Gentium} (O King of the Gentiles) & Agg 2:8; Eph 2:14, 20; Gen 2:7 & Cornerstone making two peoples one: ``come, to save man, whom Thou hast made of the dust of the earth!''\\
23 Dec & \textit{O Emmanuel} (O God-with-us) & Is 7:14, 33:22; Gen 49:10 & Our King and Lawgiver, expectation of the nations: ``come to save us, O Lord our God!''\\
\bottomrule
\end{longtable}
}

Look at the opening letters of those seven Latin titles in order:
\textbf{S}apientia, \textbf{A}donai, \textbf{R}adix, \textbf{C}lavis,
\textbf{O}riens, \textbf{R}ex, \textbf{E}mmanuel. Read backward from 23
December to 17 December, those seven initials spell two Latin words:
\textbf{\textit{ERO CRAS}}---``Tomorrow I will be there.'' On the evening of 23
December, the seven petitions finish, and the answer stands one sleep away.

The full Latin and English texts of all seven Great Antiphons are printed in
Appendix~\ref{app:prayers}. Praying each evening's antiphon aloud before supper
or night prayer---followed by the \textit{Magnificat}---gives a family or a
solitary believer the cleanest possible rhythm for the final seven days. While
the surrounding world races through December noise, the Church slows down to
one title, one Scripture, and one cry each dusk: on 19 December, ``In that day
the root of Jesse, who standeth for an ensign of the people, him the Gentiles
shall beseech, and his sepulchre shall be glorious'' (Is 11:10); on 20
December, ``And I will lay the key of the house of David upon his shoulder''
(Is 22:22); on 21 December, ``Through the bowels of the mercy of our God, in
which the Orient from on high hath visited us'' (Luke 1:78). In the longest
Advents---when the Fourth Week reaches Thursday (two of the seven calendar
shapes) and Friday (one of the seven) before the Vigil, since impeded ferial
lessons are omitted under the 1960 rubrics (RGBR n.~213)---Matins closes the
book of Isaiah with chapter 64 (``O that thou wouldst rend the heavens, and
wouldst come down'') and chapter 66; the home reading course in
Appendix~\ref{app:daily-course} keeps both chapters on 21 and 22 December every
year, until every prophet falls silent on the morning of 24 December.
